\documentclass{article}

\usepackage[utf8]{inputenc}
\usepackage[T2A]{fontenc}
\usepackage[russian]{babel}
\usepackage{amsmath}

\begin{document}
\small

\section*{Task 1.4}

\subsection*{1) Ускорили 10\% кода в 10000 раз}

Доля части которую ускорили $f=0.1$, а само ускорение $k=10000$.
Тогда:
\[
S=\frac{1}{(1-f)+\frac{f}{k}}=\frac{1}{0.9+\frac{0.1}{10000}}=\frac{1}{0.90001}\approx 1.1111
\]
Ускорение примерно 11\%.

\subsection*{2) Сколько CPU нужно для 100 минут}

На 1 CPU: $A=50$ минут (последовательно), $B=250$ минут (параллельно), всего $300$.
Пусть T(N) - время исполнения, N - количество процессоров. Тогда на $N$ CPU время:
\[
T(N)=A+\frac{B}{N}=50+\frac{250}{N}
\]
Нужно $T(N)\le 100$:
\[
50+\frac{250}{N}\le 100 \Rightarrow \frac{250}{N}\le 50 \Rightarrow N\ge 5
\]
нужно как минимум нужно $5$ CPU.

\subsection*{3) Пример: 10x ускорение только при 100 CPU}

Пусть задача состоит из 2 частей - последовательной и параллельной:
\[
A=10\text{ мин (последовательная)},\quad B=100\text{ мин (параллельная)}
\]
На 1 CPU:
\[
T(1)=A+B=110\text{ мин}
\]
На 100 CPU:
\[
T(100)=A+\frac{B}{100}=10+\frac{100}{100}=11\text{ мин}
\]
Ускорение:
\[
S=\frac{T(1)}{T(100)}=\frac{110}{11}=10
\]
\end{document}